\section{Data and Setting}
\label{sec:data}

\subsection{The Swing Platform}
\label{sec:data_platform}

Swing is one of the largest shared e-scooter operators in South Korea, serving over 50 cities with fleets of electric scooters and electric bicycles. During 2023, Swing offered three scooter operating modes on identical hardware: TUB (turbo), STD (standard), and ECO (economy). All three modes use the same vehicle models---the software-based speed governor is the only differentiator. TUB mode imposes no effective speed governor, allowing riders to reach the vehicle's maximum mechanical speed. STD mode enforces a soft governor at approximately 25~km/h, and ECO mode enforces a stricter governor at approximately 20~km/h. Riders select their preferred mode at the start of each trip via the Swing mobile application; mode choice is unrestricted and can vary trip to trip.

This configuration creates a within-operator natural experiment. Because the hardware is identical across modes, any behavioral differences between TUB, STD, and ECO trips can be attributed to the software governor rather than to vehicle characteristics, road conditions, or infrastructure---isolating the governor as the treatment.

\subsection{Regulatory Context}
\label{sec:data_regulation}

South Korea's regulatory framework for personal mobility devices (PMDs) evolved rapidly during 2020--2023. In December 2020, an amendment to the Road Traffic Act reclassified e-scooters as a type of bicycle, permitting operation on cycling paths and lowering the minimum age to 13 with no license requirement \citep{koreaherald2020escooter}. Mounting safety concerns---fueled by a surge in e-scooter-related injuries---prompted a swift reversal: from May 2021 onward, riders must hold at least a class-2 motorized bicycle driver's license (minimum age 16), and helmet use is mandatory, though enforcement remains limited \citep{koreaherald2021license}. Throughout this period, the national speed limit for PMDs has been 25~km/h on all road types.

Our study period (February--November 2023) falls entirely within the post-May 2021 regulatory regime, providing a stable legal backdrop. Crucially, Swing's TUB mode allowed riders to exceed the 25~km/h regulatory speed limit---the software governor in STD and ECO modes enforced compliance, while TUB mode effectively left compliance to rider discretion. This creates a natural correspondence between mode choice and regulatory compliance: STD/ECO trips are mechanically compliant, while TUB trips routinely violate the speed limit (52.6\% of TUB trips exceed 25~km/h, compared with 1.1\% for STD and 0.0\% for ECO). The December 2023 TUB ban thus represents an operator-level enforcement of an existing regulatory speed limit, not a new speed restriction per se.

\subsection{Dataset}
\label{sec:data_dataset}

Our primary dataset comprises GPS-instrumented e-scooter trips from February through November 2023. Each trip record contains the full GPS trajectory (latitude, longitude, timestamp triples at $\sim$10-second intervals), a per-trip speed vector from the vehicle's onboard sensor recorded at each GPS point, trip metadata (vehicle model, operating mode, distance, duration), and anonymized user identifiers. Speed data from the onboard sensor were validated against GPS-derived speeds on a random subsample; sensor readings provide reliable per-point speed measurements suitable for behavioral analysis.

After quality filtering---removing moped-class vehicles, invalid GPS coordinates outside Korean bounds (lat 33--39$^\circ$, lon 124--132$^\circ$), trips with fewer than 5 GPS points, and implausible maximum speeds exceeding 50~km/h---the dataset contains 19,454,309 valid scooter trips from 1,001,459 unique users across 52 cities and 17 provinces. User demographics (age) were obtained from a separate trip-level dataset and linked via anonymized user identifiers; age data are available for approximately 90\% of users (median age 24, mean 28 years). To our knowledge, this is the largest dataset of e-scooter trips with per-point speed profiles analyzed in the academic literature.

For road infrastructure association, we extracted OpenStreetMap (OSM) road networks for all 50 cities using \texttt{osmnx} \citep{boeing2017osmnx} and assigned each trip's start point to the nearest road edge via a KD-tree spatial index. This yields two binary infrastructure controls: \texttt{is\_major\_road} (primary, secondary, or trunk road) and \texttt{is\_cycling\_infra} (dedicated cycleway or cycle lane). The latter captures 1.3\% of trips, evenly distributed across modes.

\subsection{The December 2023 TUB Ban}
\label{sec:data_event}

In December 2023, Swing removed TUB mode system-wide across all cities in response to regulatory pressure regarding e-scooter speed safety. The mode share shifted dramatically: TUB dropped from 54.6\% to 1.4\% of trips, with STD absorbing the majority (40.8\% $\rightarrow$ 85.8\%). ECO remained relatively stable. This policy change was exogenous to individual rider behavior---it was a platform-level decision, not a rider choice---and constitutes an ideal natural experiment for evaluating speed governor effects.

Critically, the treatment intensity varies across cities. Pre-ban TUB adoption ranged from approximately 32\% to 70\% of trips depending on the city, creating substantial cross-city variation in the ``dose'' of governor imposition. Cities with higher pre-ban TUB share experienced a larger proportional shift toward governed modes, enabling a dose-response analysis within the difference-in-differences framework.

For the causal analysis, we extend the dataset to include December 2023, constructing a city-month panel of 52 cities $\times$ 11 months (February--December 2023; 549 observations). All cross-sectional analyses use the pre-ban period only (February--November 2023).

\subsection{Descriptive Statistics}
\label{sec:data_descriptive}

Table~\ref{tab:descriptive} presents descriptive statistics for the eight primary behavioral outcomes and three manipulation-check variables, stratified by operating mode. TUB mode accounts for 52.3\% of trips (10.2 million), STD for 42.7\% (8.3 million), and ECO for 5.0\% (0.97 million). The manipulation check confirms the expected speed gradient: TUB trips average 15.7~km/h (max 24.8~km/h), STD trips 13.4~km/h (max 20.7~km/h), and ECO trips 10.7~km/h (max 16.4~km/h).

\begin{table}[htbp]
\centering
\caption{Descriptive statistics by operating mode (February--November 2023, $N = 19{,}454{,}309$ trips).}
\label{tab:descriptive}
\small
\begin{tabular}{llrrr}
\toprule
\textbf{Category} & \textbf{Variable} & \textbf{TUB} & \textbf{STD} & \textbf{ECO} \\
 & & ($n = 10.2$M) & ($n = 8.3$M) & ($n = 0.97$M) \\
\midrule
\multicolumn{5}{l}{\textit{Manipulation check (mechanical)}} \\
 & Mean speed (km/h) & 15.72 (4.25) & 13.42 (3.71) & 10.66 (3.01) \\
 & Max speed (km/h) & 24.84 (2.73) & 20.67 (2.42) & 16.39 (2.54) \\
 & P85 speed (km/h) & 23.21 (3.03) & 19.51 (3.02) & 15.31 (2.55) \\
\addlinespace
\multicolumn{5}{l}{\textit{Safety outcomes (genuinely behavioral)}} \\
 & Harsh accel (count/trip) & 0.480 (0.82) & 0.232 (0.53) & 0.020 (0.15) \\
 & Harsh decel (count/trip) & 0.675 (0.96) & 0.332 (0.63) & 0.020 (0.15) \\
 & Speed CV & 0.564 (0.17) & 0.558 (0.19) & 0.564 (0.23) \\
\addlinespace
\multicolumn{5}{l}{\textit{Demand outcomes}} \\
 & Trip distance (m) & 928 (603) & 777 (528) & 656 (453) \\
 & Travel time (s) & 354 (204) & 350 (205) & 422 (266) \\
\addlinespace
\multicolumn{5}{l}{\textit{Riding dynamics outcomes}} \\
 & Cruise fraction & 0.259 (0.14) & 0.320 (0.15) & 0.417 (0.16) \\
 & Zero-speed fraction & 0.155 (0.09) & 0.155 (0.09) & 0.167 (0.10) \\
\bottomrule
\multicolumn{5}{l}{\footnotesize Values are mean (SD). Harsh acceleration threshold: 0.5~m/s$^2$ (calibrated for $\sim$10s GPS intervals).} \\
\multicolumn{5}{l}{\footnotesize Cruise fraction: proportion of speed readings within $\pm$3~km/h of trip mean.} \\
\end{tabular}
\end{table}

Several patterns are immediately apparent. Safety outcomes show strong mode differentiation: TUB trips average 0.48 harsh acceleration and 0.68 harsh deceleration events per trip, compared to 0.23 and 0.33 for STD, and 0.02 and 0.02 for ECO---differences spanning an order of magnitude. In contrast, speed CV is nearly identical across modes (0.56--0.56), suggesting that speed variability \textit{relative to mean speed} is a rider characteristic rather than a mode effect. Demand outcomes reveal modest differences: TUB trips are approximately 19\% longer (928~m vs.~777~m for STD) but travel time is virtually identical (354~s vs.~350~s), indicating that TUB riders cover more distance in the same time. Riding dynamics show clear patterns: cruise fraction increases from 0.26 (TUB) to 0.42 (ECO), reflecting more stable speed maintenance under stronger governance, while zero-speed fraction is similar across modes (0.15--0.17).
